\section{Beam Fusion}
\label{sec:beam-fusion}

Part II asks whether the compact captures can also produce a \emph{better 3D initialization}
than the structural no-Beam baseline.  The key observation: each per-view 2D Gaussian field is
a \emph{projection} of the same unknown 3D radiance density.  3D initialization can therefore
be posed as reconstruction-from-projections --- tomography --- instead of discrete
correspondence search.  Gaussians make this tractable in closed form: the back-projection of a
2D Gaussian is an analytic elongated 3D Gaussian, and products/fusions of Gaussians remain
Gaussian.  Beam Fusion is deterministic, RGB-free, correspondence-free (association emerges
from density overlap, not discrete matching), and CPU-first.

\subsection{Back-projection: from 2D splats to 3D beams}
\label{sec:beams}

Fix view $v$ with camera center $\bm{o}_v$ and a fitted component $(\bm{m}_j, C_j, a_j,
\bm{c}_j)$.  Its center pixel defines a world ray direction $\bm{d}_{vj}$.  At a hypothesised
camera depth $t$, the projective Jacobian of the pinhole map at the point
$\bm{p} = t\,\bm{r}_{vj}$ (camera frame) is
\begin{equation}
  J(\bm{p}) \;=\;
  \begin{pmatrix}
    f_x / z & 0 & -f_x\, x / z^2\\
    0 & f_y / z & -f_y\, y / z^2
  \end{pmatrix},
  \qquad \bm{p} = (x, y, z)\T .
  \label{eq:jacobian}
\end{equation}
Let $B = (\bm{b}_1, \bm{b}_2, \bm{r})$ be an orthonormal camera-frame basis whose third column
is the ray direction.  The fitted 2D covariance is lifted to the ray-orthogonal tangent plane
through the inverse of the projected tangent Jacobian,
\begin{equation}
  \Sigma_{\perp}
  \;=\;
  \big(J B_{:,1:2}\big)^{-1}\, C_j\, \big(J B_{:,1:2}\big)^{-\mathsf{T}},
  \label{eq:tangent-lift}
\end{equation}
which is exact for the center ray: it answers ``which tangent-plane covariance would have
projected to the observed $C_j$?''.  Along the ray the component's depth is unknown within
the working depth range $[t_{\mathrm{lo}}, t_{\mathrm{hi}}]$ obtained by intersecting the ray
with the scene bounding box (scaled camera-cloud extent, clamped to a near plane); the beam
assigns it a long variance $L^2$ with half-length
$L = \tfrac12 (t_{\mathrm{hi}} - t_{\mathrm{lo}})$.  The \emph{beam} is the 3D Gaussian with
precision
\begin{equation}
  \Lambda_{vj}(t)
  \;=\;
  R_v\T \, B \begin{pmatrix} \Sigma_{\perp}^{-1} & 0 \\ 0 & 1/L^2 \end{pmatrix} B\T R_v
  \label{eq:beam-precision}
\end{equation}
in world coordinates: certain transversally exactly as the image measured, agnostic along the
ray.  No voxel grid is ever built.

\subsection{Pair seeding in closed form}
\label{sec:seeding}

For components $i$ in view $u$ and $j$ in view $v$, the closest points between their center
rays have closed-form depths $(s, t)$ (the standard two-ray least-squares solution; parallel
rays are rejected).  A candidate 3D location is the midpoint of the closest-point segment.
Candidates are gated by three tests, all analytic:
\begin{equation}
  \underbrace{\frac{\norm{\bm{p}_u(s) - \bm{p}_v(t)}^2}
       {\sigma_u^2(s) + \sigma_v^2(t)} \le 3^2}_{\text{transverse gate}},
  \qquad
  \underbrace{\norm{\bm{c}_i - \bm{c}_j} \le 0.35}_{\text{color gate}},
  \qquad
  \underbrace{s \in [t^u_{\mathrm{lo}}, t^u_{\mathrm{hi}}],\;
              t \in [t^v_{\mathrm{lo}}, t^v_{\mathrm{hi}}]}_{\text{depth validity}},
  \label{eq:gates}
\end{equation}
where $\sigma_u(s)$ is component $i$'s mean pixel footprint scaled to world units at depth $s$
--- i.e.\ the ray separation is measured \emph{against both beams' transverse footprints}, so
large splats tolerate larger separations.  Each surviving seed receives the first-order
Gaussian product weight
\begin{equation}
  w \;=\; \tfrac12 (a_i + a_j)\,
          \exp\!\big({-\tfrac12\, g}\big)\,
          \exp\!\big({-\norm{\bm{c}_i - \bm{c}_j}^2 / (2\sigma_c^2)}\big),
  \qquad \sigma_c = 0.25,
  \label{eq:seed-weight}
\end{equation}
with $g$ the transverse gate value: negligible cross terms are never materialized.  View
pairs are processed in increasing camera-distance order; the production path streams all
pairs through a bounded 3D voxel table that retains only the strongest seed per voxel
(deterministic tie-breaking by lexicographic rank), so seed storage is capped independent of
scene size.

\subsection{Fusion by covariance intersection}
\label{sec:ci-fusion}

Contributing beams of one physical splat are \emph{correlated} observations --- they look at
the same surface through overlapping information.  The naive Gaussian product
$\Lambda = \sum_k \Lambda_k$ is therefore rejected by design: it double-counts every axis
observed by more than one view (two orthogonal views of an isotropic blob would fuse to half
its variance).  Beam Fusion uses equal-weight covariance intersection
\citep{julier1997_ci}:
\begin{equation}
  \Lambda \;=\; \frac{1}{K} \sum_{k=1}^{K} \Lambda_k,
  \qquad
  \bm{m} \;=\; \Lambda^{-1} \Big( \frac{1}{K} \sum_{k=1}^{K} \Lambda_k\, \bm{m}_k \Big),
  \label{eq:ci}
\end{equation}
where $\bm{m}_k$ is beam $k$'s world-space mean at its implied depth.  The common $1/K$
leaves the fused mean identical to the product rule but keeps the covariance conservative:
CI is exact on directions all views share and never overconfident elsewhere, which is the
standard consistency guarantee under unknown correlation.  Consequently Beam Fusion's main
failure risk is \emph{false contributor association}, not an algebraic missing factor ---
which is why contributor counts, gates, and lineage are first-class outputs
(\cref{sec:reduction}).

\subsection{Why at least three views: an observability result}
\label{sec:observability}

\begin{proposition}[Two views leave one covariance coordinate unobservable]
\label{prop:rank5}
Fix a 3D Gaussian with triangulated mean and EWA covariance projection
$C_v = A_v \Sigma A_v\T$ into views $v = 1, 2$ with world ray directions
$\bm{d}_1, \bm{d}_2$ through the mean.  The stacked linear map
$\Sigma \mapsto (C_1, C_2)$ on $\Sym(3) \cong \R^6$ has rank five for generic camera pairs,
and the symmetric matrix $\bm{d}_1 \bm{d}_2\T + \bm{d}_2 \bm{d}_1\T$ spans a shared null
direction.  Three generic views raise the rank to six.
\end{proposition}

\emph{Proof sketch.}  Each projection Jacobian annihilates its own viewing ray,
$A_v \bm{d}_v = 0$.  For the null candidate $\Delta = \bm{d}_1 \bm{d}_2\T + \bm{d}_2
\bm{d}_1\T$: $A_1 \Delta A_1\T = (A_1\bm{d}_1)(A_1\bm{d}_2)\T + (A_1\bm{d}_2)(A_1\bm{d}_1)\T
= 0$, and symmetrically for view 2, so $\Delta$ is invisible to both views; a dimension count
over $\Sym(3)$ gives generic rank five with exactly this shared mode, and a third generic ray
breaks it.  The full statement is verified symbolically and numerically in the repository
test suite.\evidence{ara/logic/claims.md C23;
tests/test_inverse_projection_fiber.py}\hfill$\square$

Two-view fusions would therefore carry an arbitrary along-baseline covariance component ---
not merely a noisy one.  Beam Fusion requires $K \ge 3$ contributing views per output
component (\texttt{min\_views}=3), turning an algebraic blind spot into a hard admission
rule.

\subsection{Greedy fold-in, reduction, and lineage}
\label{sec:reduction}

Each surviving seed projects its provisional 3D midpoint into every remaining view and
absorbs at most one nearest fitted component per view, subject to a $3\sigma$ pixel gate
(relative to the candidate component's own footprint), the color gate of \cref{eq:gates},
and a positive-depth test.  Components with fewer than three contributors are discarded
(\cref{prop:rank5}).  Each retained component is CI-fused (\cref{eq:ci}) over all its
contributors, weighted $w \cdot K$, and reduced in two exact steps: (i) contributor-signature
dedupe (identical view/component sets collapse to the strongest instance), and (ii) per-voxel
non-maximum suppression by weight --- \emph{selection}, not moment matching, so fused
covariances survive unblurred.  Covariance eigenvalues are clamped to
$[\sigma_{\min}^2, \sigma_{\max}^2]$ world bounds; colors initialize as amplitude-weighted
contributor means; opacity initializes at a constant $0.1$ (\cref{sec:carrier} explains why
no opacity can be inferred from the captures).  The output is a standard 3D Gaussian set
\emph{plus complete contributor lineage} (a CSR map from every 3D component to its
supporting 2D splats and implied depths) and per-view unmatched-splat lists --- the
correspondence by-product that downstream refinement, diagnostics, and future densification
consume.  Diagnostics recorded per run: contributor-count histogram, evaluated ray pairs,
gate survival counts, seed-voxel occupancy, and stage timings.\evidence{src/rtgs/lift/beam\_fusion.py}

\begin{figure}[t]
  \centering
  \figplaceholder{6.5cm}{%
    \textbf{Figure 7 --- Beam Fusion geometry.}  Two-panel schematic.  (a) One 2D Gaussian in
    two views back-projected as two elongated beams (transverse footprint from
    \cref{eq:tangent-lift}, long axis $L$ from the depth interval); their closest-ray
    midpoint, the transverse gate as shaded tolerance tubes, and the CI-fused 3D Gaussian at
    the intersection --- visually contrast CI (conservative ellipsoid) with the naive
    product (overconfident small ellipsoid).  (b) Fold-in: the provisional point projecting
    into a third view, the $3\sigma$ pixel gate around the nearest 2D component, and the
    resulting 3-contributor lineage.  Suggested: TikZ 3D schematic; exact ellipse shapes can
    be computed from a toy two-camera configuration.}
  \caption{\redcaption{Beam back-projection, gating, and covariance-intersection fusion.
    Schematic to be drawn.}}
  \label{fig:beam-geometry}
\end{figure}

\begin{figure}[t]
  \centering
  \figplaceholder{6cm}{%
    \textbf{Figure 8 --- initialization comparison (visual).}  Renders (front + side/top
    view) of the initial 3D Gaussian sets on one scene: Beam Fusion, balanced compact carve,
    calibration-bounded random, and (if available) COLMAP sparse init; same count budget.
    Side/top views expose along-ray elongation and floaters that front views hide.  Include
    contributor-count coloring for the Beam panel (3, 4, $\ge$5 views).  Data: existing
    development artifacts under \repopath{runs/} or a rerun of the three initializers on the
    development scene; render with \repopath{rtgs view} / preview tooling.}
  \caption{\redcaption{Initial 3D Gaussians per initializer, front and side views.  To be
    produced from checked-in initializers.}}
  \label{fig:init-visual}
\end{figure}
